% Termo de Compromisso de Autoria e Originalidade (plágio) - UFAPE
% {} deixa o campo em branco.
\documentclass[11pt,a4paper]{article}

% BLOCO DE DADOS: ALTERE APENAS OS CAMPOS ABAIXO
% Cabeçalho institucional, terceira linha
% Nome do curso, ex.: Bacharelado em Ciência da Computação
\newcommand{\FCursoHeader}{Bacharelado em Ciência da Computação}

% Identificação do trabalho, vai no corpo da declaração
% Nome completo do(a) discente, ex.: Maria da Silva
\newcommand{\FDiscente}{}
% Nome do curso, ex.: Ciência da Computação
\newcommand{\FCursoTrabalho}{Ciência da Computação}
% Tipo por extenso, ex.: Trabalho de Conclusão de Curso
\newcommand{\FTipoTCC}{Trabalho de Conclusão de Curso}
% Documento por extenso, ex.: Monografia de Conclusão de Curso
\newcommand{\FTipoDoc}{Monografia de Conclusão de Curso}
% Início do período, ex.: março de 2025
\newcommand{\FPeriodoDe}{maio de 2026}
% Fim do período, ex.: dezembro de 2025
\newcommand{\FPeriodoAte}{setembro de 2026}
% Local de execução, ex.: UFAPE -- Avenida Bom Pastor, s/n. -- Boa Vista, Garanhuns/PE
\newcommand{\FLocalExec}{UFAPE -- Avenida Bom Pastor, s/n. -- Boa Vista, Garanhuns/PE}

% Local e data, vão na última linha
% Cidade, ex.: Garanhuns
\newcommand{\FCidade}{Garanhuns}
% Dia, em número, ex.: 10
\newcommand{\FDia}{}
% Mês, por extenso, ex.: Janeiro
\newcommand{\FMes}{}
% Ano, em número, ex.: 2026
\newcommand{\FAno}{}

\usepackage[T1]{fontenc}
\usepackage{mathptmx}
\usepackage[brazilian]{babel}
\usepackage{graphicx}
\usepackage{setspace}
\usepackage{etoolbox}
\hyphenation{ma-tri-cu-la-do Con-clu-sao}
\usepackage[a4paper,left=3cm,right=3cm,top=1.6cm,bottom=1.6cm]{geometry}
\usepackage[bookmarksnumbered=true,psdextra,pdfencoding=auto]{hyperref}
\hypersetup{hidelinks}

\pagestyle{empty}
\setlength{\parindent}{0pt}
\setlength{\parskip}{0pt}
\setlength{\emergencystretch}{1em}

\newcommand{\gap}[2]{\vspace{#2mm}}

\newcommand{\LAss}{7.2cm}

\newcommand{\assvazia}[1]{%
  \begin{minipage}[t]{\LAss}\centering
    \rule{\LAss}{0.4pt}\\
    \textbf{\small #1}
  \end{minipage}}

\begin{document}

{\setlength{\unitlength}{1mm}%
\begin{picture}(0,0)
  \put(-17.3,-10.4){\includegraphics[height=2.20cm]{figuras/brasao-ufape.pdf}}%
  \put(144,-10.4){\includegraphics[height=2.20cm]{figuras/marca-mec.pdf}}%
\end{picture}}%
\noindent\makebox[\textwidth][c]{%
  \begin{minipage}{14.2cm}\centering\bfseries\fontsize{11}{18.1}\selectfont
    Ministério da Educação\\
    Universidade Federal do Agreste de Pernambuco\\
    \ifdefempty{\FCursoHeader}{}{\FCursoHeader}%
  \end{minipage}%
}

\gap{cab-titulo}{3.5}
\noindent\makebox[\textwidth][c]{%
  \begin{minipage}{15cm}\centering\bfseries\fontsize{12}{15}\selectfont
    Termo de Compromisso de Autoria e Originalidade (Plágio)
  \end{minipage}%
}

\gap{titulo-1}{4.5}

\pdfbookmark[1]{1 - Declaracao de autenticidade e originalidade}{sec1}
\noindent\textbf{1- Declaração de autenticidade e originalidade:}

\gap{1}{0.5}
\begin{spacing}{1.22}
O(a) discente\ifdefempty{\FDiscente}{}{ \textbf{\FDiscente}}, abaixo-assinado, regularmente matriculado(a) no \FTipoTCC\ do curso de \FCursoTrabalho, no período de \FPeriodoDe\ a \FPeriodoAte, realizado nas dependências do(a) \FLocalExec, declara que o conteúdo do(a) \FTipoDoc\ é autêntico, original e de sua exclusiva autoria e declara ainda ciência de que, na existência de plágio, estará sujeito(a) às penalidades cabíveis e previstas na lei de Direitos Autorais.
\end{spacing}

\gap{1-2}{3.0}

\pdfbookmark[1]{2 - Das responsabilidades}{sec2}
\noindent\textbf{2- Das responsabilidades:}

\gap{2}{0.5}
\begin{spacing}{1.22}
A prática de plágio importa violação aos direitos de autor, nos termos da Lei n.\ 9.610, de 19 de fevereiro de 1998, e sujeita o(a) discente às seguintes consequências:

a) \textbf{âmbito administrativo:} a improbidade na execução dos trabalhos acadêmicos constitui infração disciplinar, passível das penalidades de advertência, suspensão, reparação de danos ao patrimônio e desligamento, conforme o Manual de Conduta e Convivência Discente da UFAPE (Resolução n.\ 009/2023-CONSUNI);

b) \textbf{âmbito penal:} conforme o art.\ 184 do Código Penal (Decreto-Lei n.\ 2.848, de 7 de dezembro de 1940), a violação a direito de autor é punida com detenção de 3 (três) meses a 1 (um) ano, ou multa; o seu parágrafo terceiro agrava a pena para reclusão de 2 (dois) a 4 (quatro) anos, e multa, quando a obra for disponibilizada ao público por sistema que permita ao usuário selecioná-la para recebê-la em tempo e lugar por este determinados.
\end{spacing}

\gap{2-3}{3.0}

\pdfbookmark[1]{3 - Da ciencia}{sec3}
\noindent\textbf{3- Da ciência:}

\gap{3}{0.5}
\begin{spacing}{1.22}
Declaro ainda estar ciente de que a verificação da originalidade deste trabalho por meio de programa de detecção de semelhanças é de extrema relevância, e que a simples apresentação do manuscrito não exime o(a) discente de responder por eventuais similaridades entre o seu trabalho e o de terceiros.
\end{spacing}

\gap{texto-ass}{12.3}

\pdfbookmark[1]{Assinaturas}{secAss}
\noindent\begin{minipage}[t]{\textwidth}
\assvazia{Assinatura do(a) discente}%
\hfill
\assvazia{Assinatura do(a) orientador(a)}

\gap{ass-data}{11.3}
\noindent\makebox[\textwidth][c]{%
  \FCidade, \FDia\ de \FMes\ de \FAno.}
\end{minipage}

\end{document}
